\section{TCP baseline bổ trợ}

Các phép đo TCP trước đợt xác nhận được giữ làm baseline mạng. Chúng không được
trộn vào năm block RTT hoặc ba block ứng dụng mới vì khác thời điểm và thiết kế.

\begin{table}[H]
\centering
\caption{Tóm tắt TCP baseline lịch sử}
\label{tab:tcp-baseline-summary}
\small
\begin{tabularx}{\textwidth}{>{\raggedright\arraybackslash}p{0.28\textwidth}p{0.17\textwidth}p{0.17\textwidth}X}
\toprule
\textbf{Metric} & \textbf{Peering} & \textbf{TGW} & \textbf{Diễn giải} \\
\midrule
Single-flow receiver goodput & 4,782804 Gbps & 4,608265 Gbps & Peering cao hơn 3,7875\% \\
Fan-in tổng xấp xỉ & 11,611009 Gbps & 11,590520 Gbps & Chênh 0,17646\% \\
CPU busy receiver B & 59,1802\% & 56,8751\% & Không loại trừ bottleneck một core/ENA \\
Jain index fan-in & 0,999133 & 0,984494 & Mô tả chia sẻ trong phiên \\
\bottomrule
\end{tabularx}
\end{table}

Single-flow cho thấy Peering có receiver goodput cao hơn trong hai phiên đã
đo. Tuy nhiên retransmission khác mạnh giữa hai phiên: Peering ghi
48.253--57.185, TGW ghi 0--671. Vì vậy không quy toàn bộ chênh lệch goodput cho
network service.

Trong fan-in, tổng goodput của hai mode gần nhau nhưng receiver B đồng thời ghi
delta lớn ở \nolinkurl{bw_in_allowance_exceeded} và
\nolinkurl{pps_allowance_exceeded}. Ba sender cùng hội tụ vào B nên endpoint
allowance có thể che khác biệt của network path. Kết quả fan-in không phải
equivalence test và không chứng minh hai dịch vụ có throughput bằng nhau.

Baseline TCP trả lời khía cạnh throughput; đợt mới trả lời khả năng lặp lại của
RTT và cơ chế tail latency ứng dụng. Hai lớp bằng chứng bổ sung cho nhau, không
được dùng thay thế lẫn nhau.
